% Zone „kennst du schon" – Potenz- und Exponentialfunktionen, Kl. 10
\setcounter{aufgabe}{0}
\hypertarget{zone}{}\einheitenkopf*{Kennst du schon}
\zweigzeile{Das kennst du schon}

\begin{aufgabe}{Ich kann Dezimalzahlen malnehmen und das Ergebnis sinnvoll runden. Rechne im Kopf.}
\begin{teilezwei}
\tz $0{,}5 \cdot 8 =$ \leerfeld & \tz $3 \cdot 1{,}1 =$ \leerfeld \\
\tz $250 \cdot 1{,}07 =$ \leerfeld & \\
\end{teilezwei}
Rechne mit dem Taschenrechner und runde wie angegeben.
\begin{teile}
\teil $48{,}356\,€$ auf Cent: \leerfeld[€]
\teil $1\,234 \cdot 1{,}15$ Bakterien, auf ganze Bakterien: \leerfeld
\teil $37 \cdot 1{,}15 \cdot 1{,}15$, erst am Ende auf eine Nachkommastelle: \leerfeld
\end{teile}
\end{aufgabe}

\begin{aufgabe}{Ich kann einen Wert um einen Prozentsatz erhöhen oder senken, indem ich mit einer Dezimalzahl malnehme. Gib an, mit welcher Zahl du malnimmst.}
\begin{teilezwei}
\tz Erhöhung um $10\,\%$: mal \leerfeld & \tz Senkung um $50\,\%$: mal \leerfeld \\
\tz Erhöhung um $7\,\%$: mal \leerfeld & \tz Senkung um $15\,\%$: mal \leerfeld \\
\tz Erhöhung um $2{,}5\,\%$: mal \leerfeld & \\
\end{teilezwei}
Berechne den neuen Preis.
\begin{teile}
\teil $300\,€$ werden um $12\,\%$ erhöht. \quad neuer Preis: \leerfeld[€]
\teil $860\,€$ werden um $10\,\%$ gesenkt. \quad neuer Preis: \leerfeld[€]
\end{teile}
\end{aufgabe}

\begin{aufgabe}{Ich finde den Fehler beim Erhöhen um Prozent. Mia soll $400\,€$ um $2{,}5\,\%$ erhöhen. Sie rechnet so:}
\rechnung{400 \cdot 1{,}25 &= 500 && (€)}
\begin{teile}
\teil Finde den Fehler, benenne ihn und korrigiere die Rechnung.
\teil Erhöhe selbst $800\,€$ um $3{,}5\,\%$. \quad neuer Preis: \leerfeld[€]
\end{teile}
\end{aufgabe}

\begin{aufgabe}{Ich kann Punkte in einem Koordinatensystem mit großen Schritten ablesen. Lies die Koordinaten der Punkte ab.}

\begin{ksys}[xmin=0,xmax=8,ymin=0,ymax=90,ystep=10,ablesen]
\punkt{1}{50}{A}
\punkt{3}{40}{B}
\punkt{6}{70}{C}
\punkt{7}{10}{D}
\end{ksys}

\begin{teilezwei}
\tz \punktfeldn{A} & \tz \punktfeldn{B} \\
\tz \punktfeldn{C} & \tz \punktfeldn{D} \\
\end{teilezwei}
\begin{teile}
\teil Schreibe die vier Punkte als Wertetabelle.
\end{teile}
\sachtabelle{lcccc}{$x$ & \leerzelle & \leerzelle & \leerzelle & \leerzelle}{$y$ & \leerzelle & \leerzelle & \leerzelle & \leerzelle}
\end{aufgabe}

\begin{aufgabe}{Ich kann eine Reihe fortsetzen, in der jedes Mal derselbe Betrag dazukommt oder wegfällt. Setze die Reihe um eine Zahl fort.}
\begin{teilezwei}
\tz $12;\ 15;\ 18;$ \leerfeld & \tz $50;\ 70;\ 90;$ \leerfeld \\
\tz $7{,}5;\ 9;\ 10{,}5;$ \leerfeld & \tz $90;\ 83;\ 76;$ \leerfeld \\
\end{teilezwei}
Ein Becken enthält schon $40$ Liter Wasser. Jede Minute fließen $15$ Liter dazu.
\begin{teile}
\teil Wie viel Wasser ist nach $3$ Minuten im Becken? \quad \leerfeld[Liter]
\teil Du zeichnest den Graphen: waagerecht die Zeit, senkrecht die Wassermenge. Bei welchem Wert trifft er die senkrechte Achse? \quad \leerfeld[Liter]
\end{teile}
\end{aufgabe}

\begin{aufgabe}{Ich kann zwei Werte durcheinander teilen und das Ergebnis deuten. Rechne.}
\begin{teilezwei}
\tz $18 : 9 =$ \leerfeld & \tz $36 : 12 =$ \leerfeld \\
\tz $44 : 40 =$ \leerfeld & \\
\end{teilezwei}
Mit welcher Zahl muss man den ersten Wert malnehmen, damit der zweite herauskommt?
\begin{teilezwei}
\tz $50 \;\rightarrow\; 40$: \quad mal \leerfeld & \tz $400 \;\rightarrow\; 460$: \quad mal \leerfeld \\
\end{teilezwei}
\end{aufgabe}

\begin{aufgabe}{Ich kann Prozentwert und Prozentsatz berechnen. Berechne den Prozentwert.}
\begin{teilezwei}
\tz $10\,\%$ von $70\,€$: \leerfeld[€] & \tz $50\,\%$ von $36\,\text{kg}$: \leerfeld[kg] \\
\tz $7\,\%$ von $350\,€$: \leerfeld[€] & \\
\end{teilezwei}
Wie viel Prozent sind das?
\begin{teile}
\teil $14$ von $40$ Schülern: \leerfeld[\%]
\teil die Dezimalzahl $0{,}035$: \leerfeld[\%]
\teil Ein Preis steigt von $250\,€$ auf $285\,€$. Um wie viel Prozent ist er gestiegen? \quad \leerfeld[\%]
\end{teile}
\end{aufgabe}

\begin{aufgabe}{Ich kann ausrechnen, wie ein Guthaben mit Zinseszins wächst. $500\,€$ werden mit $2\,\%$ im Jahr verzinst; die Zinsen bleiben auf dem Konto und werden mitverzinst.}
\begin{teile}
\teil Guthaben nach einem Jahr: \leerfeld[€]
\teil Zinsen im ersten Jahr: \leerfeld[€]
\teil Guthaben nach zwei Jahren: \leerfeld[€]
\teil Zinsen im zweiten Jahr: \leerfeld[€]
\end{teile}
Rechne mit Anfangskapital mal Zinsfaktor hoch Anzahl der Jahre.
\begin{teile}
\teil $2\,000\,€$ werden $3$ Jahre lang mit $1{,}8\,\%$ verzinst. \quad Endkapital: \leerfeld[€]
\end{teile}
\end{aufgabe}

\begin{aufgabe}{Ich kann Potenzen mit dem Taschenrechner berechnen. Nutze die Taste \texttt{\^{}} oder $x^y$.}
\begin{teilezwei}
\tz $2^3 =$ \leerfeld & \tz $3^5 =$ \leerfeld \\
\tz $1{,}1^3 =$ \leerfeld & \tz $50 \cdot 1{,}1^3 =$ \leerfeld \\
\end{teilezwei}
Runde wie angegeben; runde erst das Endergebnis.
\begin{teile}
\teil $0{,}9^5$ auf drei Nachkommastellen: \leerfeld
\teil $1\,800 \cdot 1{,}07^{10}$ auf Cent: \leerfeld[€]
\end{teile}
\end{aufgabe}
